% !TEX root = main.tex

\section{Income Classification Model}

The classification component of this project focuses on predicting whether an individual falls into a low-income or high-income category based on demographic, employment, and socioeconomic features. This is formulated as a supervised learning problem aimed at supporting targeted marketing and resource allocation.

\subsection{Problem Definition}

The objective is to predict whether an individual's income exceeds \$50,000 using structured survey features. In addition to predictive accuracy, interpretability is an important consideration, as the model is intended to provide insight into the factors associated with income differences.

\subsection{Models Used}

Several supervised learning models were evaluated to balance predictive performance and interpretability:

\begin{itemize}
    \item \textbf{Logistic Regression (LASSO and Ridge):} Linear models used as interpretable baselines, with regularization to handle multicollinearity and high-dimensional feature spaces.
    \item \textbf{Decision Tree (Gini criterion):} A nonlinear model that captures feature interactions and provides rule-based interpretability.
    \item \textbf{XGBoost:} A gradient boosting method selected for its strong performance on structured tabular data and ability to model nonlinear relationships.
\end{itemize}

\subsection{Training Strategy}

Model performance was estimated using 5-fold stratified cross-validation to preserve class proportions across folds. Each model was trained under three weighting strategies:

\begin{itemize}
    \item Original Sample Weights
    \item Class-Balanced Weights
    \item Fully Uniform Class Weights
\end{itemize}

These strategies were used to evaluate the impact of class imbalance handling on model performance. However, all reported evaluation metrics are computed using the original survey-weighted distribution to ensure comparability.

\subsection{Evaluation Metrics}

Due to the strong class imbalance, multiple evaluation metrics are used rather than accuracy alone. ROC-AUC is used as the primary measure of discriminative ability, while F1 score, precision, recall, specificity, and negative predictive value are used to capture different aspects of classification performance.

\subsection{Results Summary}

\Cref{tab:classification-performance-across-weighting-strategies} presents classification performance for all models under the three weighting strategies.

Across all weighting schemes, XGBoost consistently achieves the strongest overall discriminative performance, particularly in terms of ROC-AUC and F1 score. Logistic regression models behave very similarly to each other and provide competitive performance across all metrics, while the decision tree model generally underperforms, especially in ROC-AUC and F1 score.

Differences between weighting strategies are primarily reflected in threshold-dependent metrics rather than ranking performance. ROC-AUC remains highly stable across all models and weighting schemes, indicating that class reweighting does not materially affect the ability of the models to rank observations by predicted risk. In contrast, class-balanced and uniform weighting consistently increase recall while substantially reducing precision and F1 score across all models, reflecting a systematic shift toward more positive classifications. Overall, the results indicate that all models separate the classes effectively in terms of ranking ability, but differ mainly in how they trade off false positives and false negatives rather than in their underlying discriminative performance.

\begin{table}[ht]
	\centering
	\small
	\setlength{\tabcolsep}{7pt}

	\caption{Classification Performance Across Weighting Strategies (mean $\pm$ standard deviation)}
	\label{tab:classification-performance-across-weighting-strategies}

	\sisetup{
		round-mode = places,
		round-precision = 4,
		detect-all = true,
		detect-weight = true,
		detect-inline-weight = math
	}

	\begin{tabular}{
		l
		S[table-format=1.4] @{\,\(\pm\)\,} S[table-format=1.4]
		S[table-format=1.4] @{\,\(\pm\)\,} S[table-format=1.4]
		S[table-format=1.4] @{\,\(\pm\)\,} S[table-format=1.4]
		S[table-format=1.4] @{\,\(\pm\)\,} S[table-format=1.4]
	}
		\toprule

		\textbf{Model}
			& \multicolumn{2}{c}{\textbf{ROC-AUC}}
			& \multicolumn{2}{c}{\textbf{F1}}
			& \multicolumn{2}{c}{\textbf{Precision}}
			& \multicolumn{2}{c}{\textbf{Recall}} \\
		\midrule

		\addlinespace
		\multicolumn{9}{c}{\textbf{Original Sample Weights}} \\
		\addlinespace

		LASSO Logistic Regression
			& 0.9428 & 0.0022
			& 0.4998 & 0.0085
			& 0.8379 & 0.0053
			& 0.6867 & 0.0041 \\

		Ridge Logistic Regression
			& 0.9428 & 0.0022
			& 0.5015 & 0.0086
			& 0.8368 & 0.0058
			& 0.6879 & 0.0041 \\

		Decision Tree
			& 0.7163 & 0.0035
			& 0.4646 & 0.0058
			& 0.7117 & 0.0051
			& \textbf{0.7160} & 0.0032 \\

		XGBoost
			& \textbf{0.9510} & 0.0018
			& \textbf{0.5504} & 0.0143
			& \textbf{0.8587} & 0.0076
			& 0.7118 & 0.0071 \\

		\midrule
		\addlinespace
		\multicolumn{9}{c}{\textbf{Class-Balanced Weights}} \\
		\addlinespace

		LASSO Logistic Regression
			& 0.9436 & 0.0022
			& 0.4384 & 0.0056
			& 0.6410 & 0.0024
			& 0.8722 & 0.0038 \\

		Ridge Logistic Regression
			& 0.9435 & 0.0021
			& 0.4388 & 0.0052
			& 0.6412 & 0.0022
			& 0.8723 & 0.0036 \\

		Decision Tree
			& 0.7073 & 0.0026
			& 0.4457 & 0.0050
			& \textbf{0.7002} & 0.0042
			& 0.7070 & 0.0026 \\

		XGBoost
			& \textbf{0.9506} & 0.0018
			& \textbf{0.4580} & 0.0072
			& 0.6498 & 0.0032
			& \textbf{0.8785} & 0.0041 \\

		\midrule
		\addlinespace
		\multicolumn{9}{c}{\textbf{Uniform Class Weights}} \\
		\addlinespace

		LASSO Logistic Regression
			& 0.9438 & 0.0022
			& 0.4357 & 0.0046
			& 0.6397 & 0.0020
			& 0.8724 & 0.0034 \\

		Ridge Logistic Regression
			& 0.9437 & 0.0020
			& 0.4363 & 0.0047
			& 0.6400 & 0.0020
			& 0.8725 & 0.0033 \\

		Decision Tree
			& 0.7098 & 0.0054
			& 0.4495 & 0.0097
			& \textbf{0.7018} & 0.0061
			& 0.7094 & 0.0054 \\

		XGBoost
			& \textbf{0.9507} & 0.0017
			& \textbf{0.4556} & 0.0067
			& 0.6486 & 0.0030
			& \textbf{0.8794} & 0.0032 \\

		\bottomrule
	\end{tabular}
\end{table}

\subsection{Model Interpretability and Key Drivers of Income}

Interpretability analysis was conducted using logistic regression coefficients, decision tree feature importance, and SHAP values for XGBoost. Despite differences in modeling approaches, a consistent set of predictive drivers emerges across all models. Higher education levels (particularly graduate and professional degrees), increased work intensity (weeks worked per year), and high-status occupations are consistently associated with higher income. In contrast, lower education levels, non-filing tax status, and being female are associated with lower predicted income. Financial variables such as capital gains, dividends, and capital losses also play a significant role, particularly in tree-based models. LASSO and Ridge logistic regression highlight education and occupation as the strongest predictors. The decision tree places greater emphasis on financial variables, while XGBoost captures a broader combination of demographic, financial, and occupational effects through nonlinear interactions. Overall, the consistency across models suggests a stable underlying structure in the data, where income is primarily driven by education, labor market engagement, and financial activity, with additional nonlinear effects captured by more flexible models.

